% Options for packages loaded elsewhere
\PassOptionsToPackage{unicode}{hyperref}
\PassOptionsToPackage{hyphens}{url}
\PassOptionsToPackage{dvipsnames,svgnames,x11names}{xcolor}
%
\documentclass[
  11pt,
]{article}
\usepackage{amsmath,amssymb}
\usepackage{setspace}
\usepackage{iftex}
\ifPDFTeX
  \usepackage[T1]{fontenc}
  \usepackage[utf8]{inputenc}
  \usepackage{textcomp} % provide euro and other symbols
\else % if luatex or xetex
  \usepackage{unicode-math} % this also loads fontspec
  \defaultfontfeatures{Scale=MatchLowercase}
  \defaultfontfeatures[\rmfamily]{Ligatures=TeX,Scale=1}
\fi
\usepackage{lmodern}
\ifPDFTeX\else
  % xetex/luatex font selection
\fi
% Use upquote if available, for straight quotes in verbatim environments
\IfFileExists{upquote.sty}{\usepackage{upquote}}{}
\IfFileExists{microtype.sty}{% use microtype if available
  \usepackage[]{microtype}
  \UseMicrotypeSet[protrusion]{basicmath} % disable protrusion for tt fonts
}{}
\makeatletter
\@ifundefined{KOMAClassName}{% if non-KOMA class
  \IfFileExists{parskip.sty}{%
    \usepackage{parskip}
  }{% else
    \setlength{\parindent}{0pt}
    \setlength{\parskip}{6pt plus 2pt minus 1pt}}
}{% if KOMA class
  \KOMAoptions{parskip=half}}
\makeatother
\usepackage{xcolor}
\usepackage[margin=2.5cm]{geometry}
\usepackage{color}
\usepackage{fancyvrb}
\newcommand{\VerbBar}{|}
\newcommand{\VERB}{\Verb[commandchars=\\\{\}]}
\DefineVerbatimEnvironment{Highlighting}{Verbatim}{commandchars=\\\{\}}
% Add ',fontsize=\small' for more characters per line
\usepackage{framed}
\definecolor{shadecolor}{RGB}{248,248,248}
\newenvironment{Shaded}{\begin{snugshade}}{\end{snugshade}}
\newcommand{\AlertTok}[1]{\textcolor[rgb]{0.94,0.16,0.16}{#1}}
\newcommand{\AnnotationTok}[1]{\textcolor[rgb]{0.56,0.35,0.01}{\textbf{\textit{#1}}}}
\newcommand{\AttributeTok}[1]{\textcolor[rgb]{0.13,0.29,0.53}{#1}}
\newcommand{\BaseNTok}[1]{\textcolor[rgb]{0.00,0.00,0.81}{#1}}
\newcommand{\BuiltInTok}[1]{#1}
\newcommand{\CharTok}[1]{\textcolor[rgb]{0.31,0.60,0.02}{#1}}
\newcommand{\CommentTok}[1]{\textcolor[rgb]{0.56,0.35,0.01}{\textit{#1}}}
\newcommand{\CommentVarTok}[1]{\textcolor[rgb]{0.56,0.35,0.01}{\textbf{\textit{#1}}}}
\newcommand{\ConstantTok}[1]{\textcolor[rgb]{0.56,0.35,0.01}{#1}}
\newcommand{\ControlFlowTok}[1]{\textcolor[rgb]{0.13,0.29,0.53}{\textbf{#1}}}
\newcommand{\DataTypeTok}[1]{\textcolor[rgb]{0.13,0.29,0.53}{#1}}
\newcommand{\DecValTok}[1]{\textcolor[rgb]{0.00,0.00,0.81}{#1}}
\newcommand{\DocumentationTok}[1]{\textcolor[rgb]{0.56,0.35,0.01}{\textbf{\textit{#1}}}}
\newcommand{\ErrorTok}[1]{\textcolor[rgb]{0.64,0.00,0.00}{\textbf{#1}}}
\newcommand{\ExtensionTok}[1]{#1}
\newcommand{\FloatTok}[1]{\textcolor[rgb]{0.00,0.00,0.81}{#1}}
\newcommand{\FunctionTok}[1]{\textcolor[rgb]{0.13,0.29,0.53}{\textbf{#1}}}
\newcommand{\ImportTok}[1]{#1}
\newcommand{\InformationTok}[1]{\textcolor[rgb]{0.56,0.35,0.01}{\textbf{\textit{#1}}}}
\newcommand{\KeywordTok}[1]{\textcolor[rgb]{0.13,0.29,0.53}{\textbf{#1}}}
\newcommand{\NormalTok}[1]{#1}
\newcommand{\OperatorTok}[1]{\textcolor[rgb]{0.81,0.36,0.00}{\textbf{#1}}}
\newcommand{\OtherTok}[1]{\textcolor[rgb]{0.56,0.35,0.01}{#1}}
\newcommand{\PreprocessorTok}[1]{\textcolor[rgb]{0.56,0.35,0.01}{\textit{#1}}}
\newcommand{\RegionMarkerTok}[1]{#1}
\newcommand{\SpecialCharTok}[1]{\textcolor[rgb]{0.81,0.36,0.00}{\textbf{#1}}}
\newcommand{\SpecialStringTok}[1]{\textcolor[rgb]{0.31,0.60,0.02}{#1}}
\newcommand{\StringTok}[1]{\textcolor[rgb]{0.31,0.60,0.02}{#1}}
\newcommand{\VariableTok}[1]{\textcolor[rgb]{0.00,0.00,0.00}{#1}}
\newcommand{\VerbatimStringTok}[1]{\textcolor[rgb]{0.31,0.60,0.02}{#1}}
\newcommand{\WarningTok}[1]{\textcolor[rgb]{0.56,0.35,0.01}{\textbf{\textit{#1}}}}
\usepackage{longtable,booktabs,array}
\usepackage{calc} % for calculating minipage widths
% Correct order of tables after \paragraph or \subparagraph
\usepackage{etoolbox}
\makeatletter
\patchcmd\longtable{\par}{\if@noskipsec\mbox{}\fi\par}{}{}
\makeatother
% Allow footnotes in longtable head/foot
\IfFileExists{footnotehyper.sty}{\usepackage{footnotehyper}}{\usepackage{footnote}}
\makesavenoteenv{longtable}
\setlength{\emergencystretch}{3em} % prevent overfull lines
\providecommand{\tightlist}{%
  \setlength{\itemsep}{0pt}\setlength{\parskip}{0pt}}
\setcounter{secnumdepth}{-\maxdimen} % remove section numbering
\ifLuaTeX
  \usepackage{selnolig}  % disable illegal ligatures
\fi
\usepackage{bookmark}
\IfFileExists{xurl.sty}{\usepackage{xurl}}{} % add URL line breaks if available
\urlstyle{same}
\hypersetup{
  colorlinks=true,
  linkcolor={blue},
  filecolor={Maroon},
  citecolor={blue},
  urlcolor={blue},
  pdfcreator={LaTeX via pandoc}}

\title{Gostim2: A Cross-Platform, Fixed-Schedule Stimulus Delivery System\\
  for Neuroimaging and Behavioral Experiments}
\author{Christophe Pallier\\
  \small INSERM--CEA Cognitive Neuroimaging Unit, NeuroSpin, Gif-sur-Yvette, France\\
  \small \texttt{christophe@pallier.org}}
\date{}

\begin{document}

\setstretch{1.5}
\maketitle

\section{Abstract}\label{abstract}

We describe Gostim2, an open-source, cross-platform stimulus delivery
system designed for behavioral and neuroimaging experiments in which the
timing and content of all stimuli are fully specified in advance. The
software reads a standard CSV or TSV file whose rows define stimulus
onset times, durations, types, and content; no programming knowledge is
required to build an experiment. Gostim2 is written in the Go
programming language and uses the SDL3 graphics library for
hardware-accelerated rendering. Visual stimuli (images, text, colored
rectangles), audio stimuli (WAV, FLAC, MP3, OGG), video, and rapid
serial presentation streams are supported. Temporal precision rests on
two mechanisms: in the default VSYNC mode, stimulus onsets are aligned
to the monitor's vertical blanking interval, yielding frame-accurate
presentation; in Variable Refresh Rate (VRR) mode, VSYNC is disabled and
a 2 ms busy-wait loop is used to hit scheduled timestamps with
sub-millisecond accuracy on suitable hardware. All keypress events are
recorded with timestamps in a plain-text log file. Hardware trigger
output for synchronization with EEG, MEG, and fMRI recording systems is
supported via the DLP-IO8-G USB device. The software ships as a single
self-contained binary --- requiring no runtime installation --- and
provides both a graphical user interface and a command-line interface
suitable for automated or headless deployment. Gostim2 is released under
the GNU General Public License v3 and available at
https://github.com/chrplr/gostim2.

\textbf{Keywords:} stimulus presentation, timing precision,
neuroimaging, fixed schedule, fMRI, EEG, SDL3, Go

\begin{center}\rule{0.5\linewidth}{0.5pt}\end{center}

\section{Introduction}\label{introduction}

Precise control of stimulus timing is a core technical requirement in
behavioral and neuroimaging experiments. In electroencephalography
(EEG), event-related potentials are time-locked to stimulus onsets at
the level of individual milliseconds; in magnetoencephalography (MEG)
the same constraint applies; in functional MRI (fMRI), the hemodynamic
response is slower but the accurate association of stimulus events with
scanner volumes requires reliable trigger timing to within a fraction of
a TR. Experiment-control software must therefore deliver stimuli at
scheduled times with minimal and predictable jitter, record responses
with accurate timestamps, and send synchronization pulses to recording
hardware at the moment of each event.

Several software packages are widely used for this purpose. The MATLAB
Psychophysics Toolbox (Brainard, 1997; Kleiner et al., 2007) and its
Python port PsychoPy (Peirce et al., 2019) are general-purpose: they
support real-time branching, participant-contingent feedback, adaptive
procedures, and a large variety of stimulus types. Expyriment (Krause \&
Lindemann, 2014) offers a simpler Python API with a similar capability
profile. These tools are appropriate for a wide range of paradigms but
carry a corresponding complexity: experiments are written as programs,
runtime environments must be installed and maintained across laboratory
computers, and the garbage collectors of Python and MATLAB can introduce
unpredictable pauses during timing-critical sections if experiments are
not carefully constructed. E-Prime (Schneider et al., 2002) provides a
graphical design environment at the cost of being proprietary and
restricted to Windows. Web-based tools such as jsPsych (de Leeuw, 2015)
and Gorilla (Anwyl-Irvine et al., 2021) enable large-scale online data
collection but are subject to browser and OS compositor latency that
makes frame-accurate timing difficult (Bridges et al., 2020).

A substantial class of neuroimaging experiments --- particularly passive
viewing and listening paradigms, naturalistic stimulation experiments,
and many localizer tasks --- does not require any of the real-time
adaptability that makes general-purpose tools complex. In these
paradigms, the full stimulus sequence is known before the session
begins; the experiment is simply a precise playback machine. For this
class of experiment, a leaner tool is preferable: one that can be
operated without programming, deployed to a new computer without
installing a language runtime, and that focuses all of its computational
resources on temporal precision rather than on supporting a general
scripting environment.

Gostim2 is built to fill this niche. It accepts a stimulus schedule as a
plain CSV or TSV file and presents each stimulus at its scheduled time.
It does not support real-time branching or adaptive feedback; this
constraint is a deliberate design choice that enables reliable resource
pre-loading, predictable memory allocation, and a tight VSYNC-locked
rendering loop. The software compiles to a single binary with no
external dependencies and provides both a graphical user interface for
interactive laboratory use and a command-line interface for scripted or
headless operation.

\begin{center}\rule{0.5\linewidth}{0.5pt}\end{center}

\section{Software Description}\label{software-description}

\subsection{Implementation and
Dependencies}\label{implementation-and-dependencies}

Gostim2 is written in Go (version 1.25 or later) and uses the following
third-party libraries:

\begin{itemize}
\tightlist
\item
  \texttt{github.com/Zyko0/go-sdl3} --- SDL3 bindings for window
  management, 2D rendering, audio streaming, and event handling.
\item
  \texttt{github.com/BurntSushi/toml} --- TOML configuration file
  parsing.
\item
  \texttt{github.com/funatsufumiya/go-gv-video} --- decoding of
  \texttt{.gv} (LZ4-compressed RGBA) video sequences.
\item
  \texttt{go.bug.st/serial} --- serial-port communication for the
  DLP-IO8-G trigger device (no CGo required).
\end{itemize}

The SDL3 binary is not required on the target machine; it is bundled at
build time through the \texttt{go-sdl3/bin} package and loaded at
runtime, so the resulting binary is fully self-contained.

\subsection{Fixed-Schedule
Architecture}\label{fixed-schedule-architecture}

The central architectural constraint of Gostim2 is that the complete
stimulus schedule is known before execution begins. At startup, the
CSV/TSV file is parsed and validated; all referenced assets (images,
audio files, fonts) are loaded into memory. Because no disk I/O occurs
during the experiment, I/O latency cannot cause dropped frames or audio
glitches. Stimulus rendering parameters (scaled dimensions, center
coordinates) are computed from the screen configuration at load time and
cached, keeping per-frame CPU work minimal.

This architecture excludes paradigms that require real-time branching
(e.g., providing feedback contingent on the participant's response, or
adapting subsequent stimuli based on earlier responses). For paradigms
where such adaptability is needed, the companion framework
\texttt{goxpyriment} (a Go library for general-purpose experiment
programming) is more appropriate.

\subsection{Stimulus Configuration
Format}\label{stimulus-configuration-format}

An experiment is fully defined by a CSV or TSV file with four mandatory
columns:

\begin{longtable}[]{@{}
  >{\raggedright\arraybackslash}p{(\columnwidth - 2\tabcolsep) * \real{0.5000}}
  >{\raggedright\arraybackslash}p{(\columnwidth - 2\tabcolsep) * \real{0.5000}}@{}}
\toprule\noalign{}
\begin{minipage}[b]{\linewidth}\raggedright
Column
\end{minipage} & \begin{minipage}[b]{\linewidth}\raggedright
Description
\end{minipage} \\
\midrule\noalign{}
\endhead
\bottomrule\noalign{}
\endlastfoot
\texttt{onset\_time} & Scheduled onset of the stimulus, in milliseconds
from the start of the run \\
\texttt{duration} & Duration of the stimulus display or playback, in
milliseconds \\
\texttt{type} & Stimulus category: \texttt{IMAGE}, \texttt{SOUND},
\texttt{TEXT}, \texttt{BOX}, \texttt{VIDEO}, \texttt{IMAGE\_STREAM},
\texttt{TEXT\_STREAM}, \texttt{SOUND\_STREAM}, or \texttt{END} \\
\texttt{stimuli} & File path, text content, or stream specification \\
\end{longtable}

For stream types (\texttt{IMAGE\_STREAM}, \texttt{TEXT\_STREAM},
\texttt{SOUND\_STREAM}), multiple items are listed in a single
\texttt{stimuli} cell separated by \texttt{\textasciitilde{}}
characters. Per-item durations and inter-stimulus gaps can be specified
inline using \texttt{:duration:gap} suffixes, allowing complex RSVP
sequences to be expressed in a single CSV row. The \texttt{END} type
stops the experiment at a specified time; if omitted, the experiment
concludes after the last scheduled event.

Optional columns include \texttt{trigger} (a byte value sent to the
DLP-IO8-G at onset) and any number of additional metadata columns that
are passed through to the event log.

A separate TOML configuration file stores session-level parameters:
subject ID, screen resolution, window mode, display index, font, scale
factor, background color, text color, fixation color, stimuli directory,
results directory, and DLP device path. This file is created
automatically by the GUI after each session and can be passed to the CLI
with \texttt{-config\ \textless{}path\textgreater{}} to reproduce the
same setup.

\subsection{Temporal Precision}\label{temporal-precision}

Precision is the primary technical objective of Gostim2. Two operating
modes are provided, selected via command-line flags or GUI checkboxes:

\textbf{VSYNC mode} (default, \texttt{-vsync}). The rendering loop calls
\texttt{SDL\_RenderPresent}, which blocks until the monitor's vertical
blanking interval. On a 60 Hz monitor this is approximately 16.67 ms per
frame; on a 120 Hz monitor approximately 8.33 ms. To avoid missing a
target frame, the engine uses a predictive look-ahead window (nominally
half a frame duration). If the scheduled onset of the next stimulus
falls within this window before the current VBLANK, the renderer
prepares it immediately so that it appears on the next flip rather than
the one after. This mechanism ensures frame-accurate stimulus onsets
under typical operating conditions. The precision of this approach is
identical in principle to that of the Psychophysics Toolbox's
Screen(`Flip') mechanism (Brainard, 1997; Kleiner et al., 2007).

\textbf{VRR mode} (\texttt{-vrr}). VSYNC is disabled and the engine
replaces the blocking flip with a 2 ms busy-wait loop. The loop polls
the system clock and renders at the precise scheduled timestamp without
waiting for a frame boundary. On systems equipped with G-Sync or
FreeSync displays --- which present a frame immediately when signaled
rather than at a fixed interval --- this mode allows stimulus onsets to
be placed at any point in time, limited only by system clock resolution
and OS scheduling jitter. The busy-wait duration (2 ms) is a
configurable trade-off between CPU load and temporal resolution.

In both modes, Go's garbage collector could in principle pause the
timing loop. For the relatively coarse delivery loop of Gostim2 (stimuli
typically on the order of hundreds of milliseconds to seconds), GC
pauses of a few milliseconds are unlikely to cause missed frames. For
the dense RSVP stream types, GC pressure is reduced by pre-loading all
assets and avoiding allocation in the hot loop.

\subsection{Stimulus Types}\label{stimulus-types}

Gostim2 renders the following stimulus categories:

\begin{itemize}
\tightlist
\item
  \textbf{IMAGE}: A static image file displayed for the specified
  duration. Supported formats include PNG, JPEG, BMP, and any format
  handled by SDL3's image loader. Images are scaled according to the
  global scale factor and centered on screen by default.
\item
  \textbf{TEXT}: A text string rendered in the configured font and
  color.
\item
  \textbf{BOX}: A filled colored rectangle, useful for visual masks,
  blank intervals with non-default background color, or simple geometric
  stimuli.
\item
  \textbf{SOUND}: An audio file played for its natural duration or the
  specified duration. Supported formats: WAV, FLAC, MP3, OGG. All audio
  is resampled at load time to the device's native sample rate,
  eliminating resampling overhead during playback.
\item
  \textbf{VIDEO}: A video file played frame-by-frame, VSYNC-locked. The
  \texttt{.gv} (LZ4-compressed RGBA) format is natively supported; MPEG
  is also handled.
\item
  \textbf{IMAGE\_STREAM / TEXT\_STREAM}: A rapid serial sequence of
  images or text items presented within a single CSV row, with per-item
  onset asynchronies specified inline. This supports RSVP and temporal
  masking paradigms.
\item
  \textbf{SOUND\_STREAM}: A sequence of audio clips played in rapid
  succession, analogous to visual streams.
\end{itemize}

A superimposed fixation cross can be configured to appear on blank
intervals, always, or never, independently of the stimulus schedule.

\subsection{Audio}\label{audio}

Audio is handled by a custom software mixer implemented directly over
SDL3's audio stream API. Up to 16 sounds can play simultaneously; the
mixer sums their samples with saturation at 16-bit limits. All audio
files are decoded and resampled to the device's native format (default:
44100 Hz, stereo, signed 16-bit) during the pre-loading phase. Because
audio data is already in the device's format when playback begins, there
is no per-sample conversion overhead at play time and sound onset is not
delayed by decompression.

\subsection{Event Logging}\label{event-logging}

All keypress events --- including the key code, timestamp relative to
experiment start (in milliseconds), and the identity of the currently
displayed stimulus --- are recorded in a plain-text log file written to
the results directory. The filename encodes the experiment name, subject
ID, and a timestamp to prevent overwriting. Header lines record the
configuration parameters used for the session, providing a
self-contained record of each run.

\subsection{Hardware Trigger Output}\label{hardware-trigger-output}

For EEG, MEG, and fMRI synchronization, Gostim2 supports the DLP-IO8-G
USB device, which presents as a virtual serial port and controls eight
TTL output lines. The three lower lines are mapped to stimulus
categories: line 1 to IMAGE and VIDEO onsets, line 2 to SOUND onsets
(with a fixed 5 ms pulse), and line 3 to TEXT and BOX onsets. The DLP
device path is specified in the configuration; the engine opens the port
at startup and closes it cleanly on exit. An \texttt{AutoDetect}
function scans available serial ports for a DLP-IO8-G signature,
returning a no-op implementation if no device is found, so that the same
experiment file runs unmodified on systems without recording hardware.

\subsection{User Interfaces}\label{user-interfaces}

Both interfaces share the same underlying engine, configuration struct,
and event log format:

\textbf{GUI (\texttt{gostim2-gui}).} A graphical SDL3 dialog presents
labeled text fields for the configuration file path, subject ID, CSV
file, start/end splash screens, font, results directory, and DLP device
path; numeric fields for display index, font size, and scale factor;
color pickers for background, text, and fixation colors; and checkbox
toggles for VSYNC, VRR, autodetect resolution, fullscreen mode,
skip-wait, and fixation mode. Settings are persisted to
\texttt{gostim2\_config.toml} after each session and pre-filled on the
next run.

\textbf{CLI (\texttt{gostim2}).} A command-line interface accepting
flags corresponding to each configuration field. The same TOML file used
by the GUI can be loaded with
\texttt{-config\ \textless{}path\textgreater{}}, allowing the GUI to be
used for interactive setup and the CLI for automated scripted runs
(e.g., from an MRI scanner room trigger script).

\subsection{Display Modes}\label{display-modes}

Three window modes control how Gostim2 interacts with the OS display
stack:

\begin{itemize}
\tightlist
\item
  \textbf{Windowed} (\texttt{window\_mode\ =\ 0}): a resizable window at
  the specified resolution. Suitable for development and testing.
\item
  \textbf{Fullscreen desktop} (\texttt{window\_mode\ =\ 1}): a
  borderless fullscreen window at the current desktop resolution. The OS
  compositor may introduce one frame of additional buffering.
\item
  \textbf{Fullscreen exclusive} (\texttt{window\_mode\ =\ 2}): requests
  exclusive access to the display at the specified mode, bypassing the
  compositor. This mode provides the most predictable frame delivery on
  Windows and macOS.
\end{itemize}

On Linux, additional precision can be obtained by running the CLI
directly from a TTY console (with the display manager stopped), allowing
SDL3 to interface with the Direct Rendering Manager (DRM) without any
compositor overhead.

\begin{center}\rule{0.5\linewidth}{0.5pt}\end{center}

\section{Example Experiments}\label{example-experiments}

The repository includes three complete working experiments:

\begin{itemize}
\tightlist
\item
  \textbf{Jabberwocky}: An auditory language experiment presenting words
  and non-words from Lewis Carroll's poem with precise SOAs, accompanied
  by trigger pulses for EEG/MEG.
\item
  \textbf{La Cigale et la Fourmi}: A naturalistic listening experiment
  presenting a spoken French poem stimulus, illustrating the use of the
  \texttt{SOUND} type with long continuous audio.
\item
  \textbf{Visual Categories Localizer (demo)}: A visual category
  localizer adapted from Minye Zhan, presenting images from multiple
  categories (faces, objects, scenes, scrambled) at 150 ms per image in
  RSVP streams, with trigger output and a full configuration file.
\item
  \textbf{Simple Localizer (demo)}: A simple auditory localizer adapted
  from Philippe Pinel, demonstrating TEXT, IMAGE, and SOUND stimulus
  types in a single experiment.
\end{itemize}

All examples include the stimulus files, CSV schedule, and TOML
configuration needed to run them immediately after installation.

\begin{center}\rule{0.5\linewidth}{0.5pt}\end{center}

\section{Comparison with Existing
Tools}\label{comparison-with-existing-tools}

Table 1 compares Gostim2 with selected tools used for stimulus
presentation in behavioral and neuroimaging research.

\textbf{Table 1.} Comparison of selected stimulus presentation systems.

\begin{longtable}[]{@{}
  >{\raggedright\arraybackslash}p{(\columnwidth - 10\tabcolsep) * \real{0.1667}}
  >{\raggedright\arraybackslash}p{(\columnwidth - 10\tabcolsep) * \real{0.1667}}
  >{\raggedright\arraybackslash}p{(\columnwidth - 10\tabcolsep) * \real{0.1667}}
  >{\raggedright\arraybackslash}p{(\columnwidth - 10\tabcolsep) * \real{0.1667}}
  >{\raggedright\arraybackslash}p{(\columnwidth - 10\tabcolsep) * \real{0.1667}}
  >{\raggedright\arraybackslash}p{(\columnwidth - 10\tabcolsep) * \real{0.1667}}@{}}
\toprule\noalign{}
\begin{minipage}[b]{\linewidth}\raggedright
Property
\end{minipage} & \begin{minipage}[b]{\linewidth}\raggedright
Gostim2
\end{minipage} & \begin{minipage}[b]{\linewidth}\raggedright
Expyriment
\end{minipage} & \begin{minipage}[b]{\linewidth}\raggedright
PsychoPy
\end{minipage} & \begin{minipage}[b]{\linewidth}\raggedright
E-Prime
\end{minipage} & \begin{minipage}[b]{\linewidth}\raggedright
jsPsych
\end{minipage} \\
\midrule\noalign{}
\endhead
\bottomrule\noalign{}
\endlastfoot
Language & Go & Python & Python & Proprietary & JavaScript \\
Experiment specification & CSV/TSV file & Python script & Python script
/ Builder & GUI & JavaScript \\
Programming required & No & Yes & Optional & Optional & Yes \\
Deployment & Single binary & pip/conda env & Standalone or pip & Windows
installer & Browser/Node.js \\
Real-time branching & No & Yes & Yes & Yes & Yes \\
VSYNC sync & Yes (SDL3) & Yes (Pygame) & Yes (Pyglet/PTB) & Yes &
Limited (rAF) \\
Sub-frame timing (VRR) & Yes (busy-wait) & No & No & No & No \\
Custom audio mixer & Yes & No & No & No & No \\
Hardware triggers & DLP-IO8-G, parallel & Parallel, serial & Parallel,
serial & Serial & No \\
EEG/MEG/fMRI suitability & High (fixed-schedule) & Medium & Medium &
High & Low \\
License & GPL v3 & GPL v3 & GPL v3 & Proprietary & MIT \\
Reference & This paper & Krause \& Lindemann, 2014 & Peirce et al., 2019
& Schneider et al., 2002 & de Leeuw, 2015 \\
\end{longtable}

The most distinctive characteristic of Gostim2 relative to
general-purpose tools is its fixed-schedule constraint. By forgoing
real-time adaptability, the engine can pre-load all assets, avoid
dynamic memory allocation in the hot loop, and maintain a tight
VSYNC-locked rendering cycle. Researchers who need adaptive procedures,
participant feedback, or trial-by-trial parameter variation should use a
general-purpose tool. Researchers running fixed-sequence paradigms ---
the majority of fMRI localizers, MEG passive listening tasks, and many
RSVP experiments --- can describe their experiment as a spreadsheet and
rely on Gostim2 for delivery without writing any code.

The absence of a Python runtime dependency is a practical advantage in
multi-user laboratory environments: Gostim2 is distributed as a
pre-built binary for Windows (x86\_64 and ARM64), macOS (Intel and Apple
Silicon), and Linux (x86\_64 AppImage and Debian package), and requires
no installation beyond placing the binary on the target machine.

\begin{center}\rule{0.5\linewidth}{0.5pt}\end{center}

\section{Limitations and Future
Directions}\label{limitations-and-future-directions}

\textbf{Fixed-schedule constraint.} The design intentionally excludes
real-time feedback and adaptive stimulus selection. This is a feature
for the target use case but a hard limitation for paradigms requiring
contingent responses.

\textbf{No formal timing validation study.} The temporal precision
mechanisms described --- VSYNC synchronization and VRR busy-wait --- are
well-established approaches (Brainard, 1997; Kleiner et al., 2007), but
no published photodiode or oscilloscope measurements specific to Gostim2
on particular hardware configurations currently exist. Such measurements
are a priority for future work.

\textbf{Trigger mapping is fixed.} The current DLP-IO8-G trigger
assignment (line 1: images/video; line 2: sounds; line 3: text/boxes) is
hardcoded. Future versions should allow per-stimulus trigger values to
be specified directly in the CSV.

\textbf{Audio mixing is software-only.} The custom mixer operates in
software at 44100 Hz stereo signed 16-bit. It does not support
hardware-accelerated mixing or sample rates other than 44100 Hz without
pre-resampling. Audio playback latency is bounded by the SDL3 audio
buffer size and is not independently measured.

\textbf{No response-contingent triggers.} Keypresses are logged but do
not currently trigger DLP output lines, which some EEG/MEG acquisition
protocols require.

Future planned developments include per-row trigger configuration in the
CSV format, response-contingent trigger output, a formal timing
validation protocol, and support for parallel-port triggers on Linux.

\begin{center}\rule{0.5\linewidth}{0.5pt}\end{center}

\section{Open Practices Statement}\label{open-practices-statement}

The source code for Gostim2, along with pre-compiled binaries for
Windows, macOS, and Linux and complete documentation, is freely
available under the GNU General Public License v3 at
https://github.com/chrplr/gostim2. All example experiments included in
the repository (stimuli, CSV schedules, configuration files) can be run
immediately after installation to verify correct operation.

\begin{center}\rule{0.5\linewidth}{0.5pt}\end{center}

\section{Declarations}\label{declarations}

\textbf{Funding} None

\textbf{Conflicts of interest / Competing interests} The authors declare
no competing interests.

\textbf{Ethics approval} Not applicable. This manuscript describes a
software tool and contains no human subjects data.

\textbf{Consent to participate / Consent for publication} Not
applicable.

\textbf{Availability of data and materials} Not applicable.

\textbf{Code availability} The source code is hosted at
https://github.com/chrplr/gostim2 under the GNU General Public License
v3. Pre-built binaries for all supported platforms are available at
https://github.com/chrplr/gostim2/releases.

\textbf{Author Contributions}

\textbf{Christophe Pallier}: Conceptualization, Software (original
conception, architecture, all design decisions, integration and
testing), Writing --- original draft (conceptualization and framing),
Writing --- review and editing.

\textbf{Claude Sonnet 4.6 (Anthropic)}: Software (substantial
contributions to engine modules including \texttt{engine/audio.go},
\texttt{engine/gui\_setup.go}, \texttt{engine/run.go},
\texttt{engine/csv\_parser.go}, \texttt{engine/validator.go}, and
associated tests; co-development of the CLAUDE.md architectural context
document), Writing --- original draft (the revised paper text was
drafted collaboratively in an interactive session).

Note on AI authorship: consistent with current editorial norms (Nature
Portfolio, 2023; Springer Nature, 2023), Claude Sonnet 4.6 is listed
under Author Contributions rather than as a named author, as authorship
requires the capacity to assume accountability for the work. Christophe
Pallier takes full responsibility for all content, including the
AI-assisted portions.

\begin{center}\rule{0.5\linewidth}{0.5pt}\end{center}

\section{References}\label{references}

Anwyl-Irvine, A. L., Massonnié, J., Flitton, A., Kirkham, N., \&
Evershed, J. K. (2021). Gorilla in our midst: An online behavioral
experiment builder. \emph{Behavior Research Methods}, \emph{52}(1),
388--407. https://doi.org/10.3758/s13428-019-01237-x

Brainard, D. H. (1997). The Psychophysics Toolbox. \emph{Spatial
Vision}, \emph{10}(4), 433--436. https://doi.org/10.1163/156856897X00357

Bridges, D., Pitiot, A., MacAskill, M. R., \& Peirce, J. W. (2020). The
timing mega-study: Comparing a range of experiment generators, both
lab-based and online. \emph{PeerJ}, \emph{8}, e9414.
https://doi.org/10.7717/peerj.9414

de Leeuw, J. R. (2015). jsPsych: A JavaScript library for creating
behavioral experiments in a Web browser. \emph{Behavior Research
Methods}, \emph{47}(1), 1--12. https://doi.org/10.3758/s13428-014-0458-y

Donovan, A. A. A., \& Kernighan, B. W. (2016). \emph{The Go programming
language}. Addison-Wesley.

Kleiner, M., Brainard, D. H., \& Pelli, D. (2007). What's new in
Psychtoolbox-3? \emph{Perception}, \emph{36}(ECVP Abstract Supplement),
14.

Krause, F., \& Lindemann, O. (2014). Expyriment: A Python library for
cognitive and neuroscientific experiments. \emph{Behavior Research
Methods}, \emph{46}(2), 416--428.
https://doi.org/10.3758/s13428-013-0390-6

Nature Portfolio. (2023). \emph{Tools and technologies: Artificial
intelligence}.
https://www.nature.com/nature-portfolio/editorial-policies/ai

Peirce, J. W., Gray, J. R., Simpson, S., MacAskill, M. R., Höchenberger,
R., Sogo, H., Kastman, E., \& Lindeløv, J. K. (2019). PsychoPy2:
Experiments in behavior made easy. \emph{Behavior Research Methods},
\emph{51}(1), 195--203. https://doi.org/10.3758/s13428-018-01193-y

Schneider, W., Eschman, A., \& Zuccolotto, A. (2002). \emph{E-Prime
user's guide}. Psychology Software Tools.

Springer Nature. (2023). \emph{Artificial intelligence (AI) policy for
authors}. https://www.springernature.com/gp/policies/editorial-policies

Zyko0. (2024). \emph{go-sdl3: Go bindings for SDL3}. GitHub.
https://github.com/Zyko0/go-sdl3

\begin{center}\rule{0.5\linewidth}{0.5pt}\end{center}

\appendix
\section{OS Security Warnings for Unsigned
Binaries}\label{appendix-os-security-warnings-for-unsigned-binaries}

Gostim2 is distributed as a pre-compiled binary that is not signed with
a commercial code-signing certificate. Both macOS and Windows apply
security mechanisms that will warn users --- or actively block execution
--- the first time an unsigned binary downloaded from the internet is
run. This section describes the expected behavior and how to proceed.

\subsection{macOS: GateKeeper}\label{macos-gatekeeper}

Apple's GateKeeper evaluates downloaded applications against three tiers
of trust:

\begin{enumerate}
\def\labelenumi{\arabic{enumi}.}
\tightlist
\item
  \textbf{Unsigned} --- no signature at all. GateKeeper displays a
  dialog stating the application ``is damaged and can't be opened.''
  This phrasing is misleading: the binary is not corrupted; it simply
  lacks a signature.
\item
  \textbf{Ad-hoc signed} --- signed with \texttt{codesign\ -\/-sign\ -}
  at no cost, without an Apple Developer account. GateKeeper shows a
  one-time warning on first launch (``Apple cannot verify that this app
  is free of malware'') but allows the user to proceed.
\item
  \textbf{Apple-notarized} --- signed and notarized through Apple's
  notarization service, which requires a paid Apple Developer Program
  membership (\$99/year). No warning is shown.
\end{enumerate}

Gostim2 releases are distributed as ad-hoc signed \texttt{.app} bundles
(macOS) or ad-hoc signed CLI binaries, so users will see a one-time
first-launch warning rather than the alarming ``damaged'' message.

\textbf{For graphical \texttt{.app} bundles:} Right-click (or
Control-click) the \texttt{.app} file and select ``Open''. A dialog will
ask you to confirm; click ``Open''. This is required only once.

Alternatively: open System Settings → Privacy \& Security. A message
about the blocked app appears at the bottom of the page with an ``Open
Anyway'' button.

\textbf{For command-line binaries:} The quarantine extended attribute
must be removed before first use:

\begin{Shaded}
\begin{Highlighting}[]
\ExtensionTok{xattr} \AttributeTok{{-}d}\NormalTok{ com.apple.quarantine gostim2}
\FunctionTok{chmod}\NormalTok{ +x gostim2}
\end{Highlighting}
\end{Shaded}

For detailed instructions and the packaging script used to create ad-hoc
signed releases, see:
https://chrplr.github.io/note-about-macos-unsigned-apps/

\subsection{Windows: SmartScreen}\label{windows-smartscreen}

Windows SmartScreen (part of Microsoft Defender) evaluates executables
downloaded from the internet using a reputation system. An executable
that has not accumulated sufficient download history --- as is typical
for software distributed outside the Microsoft Store or without an
Extended Validation (EV) code-signing certificate --- triggers a blue
dialog: ``Windows protected your PC. Microsoft Defender SmartScreen
prevented an unrecognized app from starting.''

\textbf{To run Gostim2:} click ``More info'' in that dialog, then click
``Run anyway''. The warning appears only once per executable file.

For the NSIS-based Windows installer (\texttt{gostim2-setup.exe}), the
same dialog may appear before installation. The same ``More info'' →
``Run anyway'' procedure applies.

Obtaining an EV code-signing certificate from a commercial certificate
authority would suppress this warning automatically but involves ongoing
cost and is not currently pursued for this open-source project.

\subsection{Note for Researchers Distributing Experiments to
Participants}\label{note-for-researchers-distributing-experiments-to-participants}

If you compile Gostim2 from source and distribute the resulting binary
to participant computers, those computers will encounter the same
warnings described above. The procedures are identical: on macOS, remove
the quarantine attribute or use the right-click-Open method; on Windows,
use ``More info'' → ``Run anyway''. Informing participants or technical
staff in advance avoids confusion.

\end{document}
